\section{Metodologia}
\label{sec:methodology}

\subsection{Projeto de Pesquisa}

Este trabalho adota uma abordagem experimental controlada para comparar três bancos de dados de
séries temporais em nove padrões distintos de carga de trabalho IoT e três escalas de volume de
dados. A avaliação cobre vazão de escrita, latência por operação (média e percentil 99) e consumo
de recursos (CPU e memória) para cada combinação de banco de dados, tipo de carga de trabalho e
escala. O objetivo é caracterizar quais compromissos arquiteturais favorecem quais perfis de
aplicação, indo além de rankings agregados para explicar \textit{por que} cada banco de dados
tem o desempenho que tem.

Um objetivo central de projeto, além da medição de desempenho, é a \textit{reprodutibilidade}.
A literatura comparativa revisada neste trabalho revela uma lacuna sistêmica: dos cinco estudos
de desempenho recentes de BDSTs analisados, apenas um publicou scripts e configurações
suficientes para replicação independente~\citep{pereira2023mulletbench}. Os demais ou omitiram
configurações completamente, descreveram suas configurações em prosa vaga demais para seguir
mecanicamente, ou utilizaram cadeias de ferramentas que requerem acesso
institucional~\citep{shah2022performance,naqvi2017influxdb,lima2024avaliando,wang2023iotdb}.
Este trabalho trata a reprodutibilidade como uma saída de primeira classe: cada componente
necessário para replicar os resultados está fixado em versão e publicamente disponível junto
com este artigo.

\subsection{Seleção dos Bancos de Dados}

Os três bancos de dados foram selecionados para representar abordagens arquiteturais distintas
para armazenamento de séries temporais, cada um identificado como sistema líder na literatura
revisada.

\textbf{InfluxDB 2.7} é o padrão de referência de facto na pesquisa comparativa de BDSTs,
aparecendo em quatro dos cinco estudos de desempenho revisados~\citep{naqvi2017influxdb,
shah2022performance,lima2024avaliando,wang2023iotdb}. Ele usa um motor de armazenamento TSM
(Time-Structured Merge Tree) que organiza dados em blocos comprimidos ordenados por tempo. A
ingestão ocorre via API HTTP com protocolo de linha; consultas analíticas usam a linguagem de
expressão Flux. A ubiquidade do InfluxDB em benchmarks anteriores o torna uma âncora essencial
de comparação: qualquer novo estudo que o omita perde comparabilidade com a literatura existente.

\textbf{Apache IoTDB 1.3.0}~\cite{wang2023iotdb} emergiu como líder de desempenho nos dois
estudos comparativos mais recentes~\citep{pereira2023mulletbench,wang2023iotdb}, representando
uma mudança em relação à dominância do InfluxDB reportada em benchmarks mais antigos.
Desenvolvido na Universidade Tsinghua especificamente para IoT industrial, usa um formato de
arquivo colunar personalizado (TsFile) e um protocolo binário Thrift RPC
(SESSION\_BY\_TABLET). A ingestão de escritas é armazenada em buffer no heap gerenciado pela
JVM antes de ser descarregada para o disco como segmentos colunares ordenados, possibilitando
vazão excepcionalmente alta para fluxos de séries temporais ordenadas. Sua rápida ascensão em
benchmarks recentes, combinada com sua abordagem arquitetural distinta, tornou-o candidato
natural para este estudo.

\textbf{TimescaleDB} estende o PostgreSQL~15 com capacidades de séries temporais: particionamento
automático de hipertabela por intervalo de tempo, a função de agregação \texttt{time\_bucket()}
para subamostragem e compressão opcional por bloco. Ele mantém o motor relacional completo com
durabilidade baseada em WAL, controle de concorrência MVCC, índices B-tree e compatibilidade
completa com SQL. Para organizações com infraestrutura PostgreSQL existente, ou para aplicações
que exigem junções entre dados de séries temporais e relacionais, ele fornece um caminho para
a capacidade de séries temporais sem abandonar o ecossistema relacional. Aparece em quatro dos
seis benchmarks revisados~\citep{shah2022performance,lima2024avaliando,wang2023iotdb,grzesik2020comparative}.

O PostgreSQL simples, sem a extensão TimescaleDB, foi explicitamente excluído. O TimescaleDB
tende a superar o PostgreSQL simples neste domínio; incluir ambos duplicaria os testes.
Conclusões sobre configurações gerais do PostgreSQL, como dimensionamento de
\texttt{shared\_buffers}, são transferíveis, mas resultados que dependem de otimizações
específicas do TimescaleDB, como hipertabelas e eliminação de blocos por estatísticas de
mínimo/máximo, não se aplicam diretamente ao PostgreSQL sem a extensão.

\subsection{Ferramenta de Benchmark}
\label{sec:benchmark-tool}

O gerador de carga de trabalho é o \textit{iot-benchmark}~\cite{iotbenchmark}, desenvolvido na
Universidade Tsinghua, que vem com clientes específicos para os três bancos de dados em
avaliação. Esse suporte nativo a múltiplos bancos de dados significa que cada carga de trabalho
é expressa no protocolo correto para cada alvo: a interface Thrift SESSION\_BY\_TABLET para
IoTDB, o protocolo de linha HTTP para InfluxDB e uma conexão JDBC para TimescaleDB. Uma
interface consistente e parametrizada garante que as diferenças de desempenho medidas reflitam
os próprios bancos de dados, e não a camada de benchmarking.

Ferramentas alternativas foram avaliadas. O Time Series Benchmark Suite (TSBS) oferece um
framework maduro para múltiplos bancos de dados, mas carece de suporte nativo ao IoTDB no
momento desta escrita. O Apache JMeter, utilizado em Lima et al.~\cite{lima2024avaliando},
lida bem com sistemas baseados em HTTP, mas requer scripts personalizados para interfaces
Thrift e JDBC, introduzindo potencial sobrecarga específica da ferramenta. A interface uniforme
do iot-benchmark e a semântica de carga de trabalho específica para IoT (clientes
  concorrentes organizados em hierarquias de dispositivos e sensores, tamanhos de lote
configuráveis e geração de dados fora de ordem distribuída por Poisson) tornaram-no a
escolha mais adequada para este estudo.

A ferramenta produz duas matrizes de saída estruturadas que são analisadas pela camada de
orquestração. A Matriz de Resultado reporta operações totais, tempo decorrido e vazão agregada
(pts/s). A Matriz de Latência reporta latências média, mínima, máxima e por percentil por tipo
de operação, fornecendo os dados de comportamento de cauda por operação necessários para
avaliar cargas de trabalho IoT em tempo real.

\subsection{Projeto das Cargas de Trabalho}
\label{sec:workload}

O benchmark cobre nove tipos de carga de trabalho em duas categorias, projetados para abranger
o espectro de padrões encontrados em sistemas IoT em produção.

\subsubsection{Testes de Escrita}

Quatro testes de escrita exercitam padrões de ingestão distintos:

\begin{description}
  \item[\texttt{write}] Ingestão sequencial de linha de base com carimbos de tempo ordenados e
    o tamanho de lote padrão de 100 linhas por requisição. Mede a vazão bruta de escrita em
    condições ideais e ordenadas.

  \item[\texttt{out-of-order}] Vinte por cento dos pontos de dados chegam com carimbos de tempo
    mais antigos do que a fronteira de ingestão atual, modelado como um atraso distribuído por
    Poisson. Simula sensores IoT que perdem conectividade e depois reproduzem leituras
    armazenadas em buffer, um padrão ubíquo em implantações industriais e de campo onde os
    sensores operam sobre redes celulares ou LPWAN não confiáveis.

  \item[\texttt{batch-small}] \texttt{BATCH\_SIZE=1}: um ponto de dado por requisição de
    escrita. Isola a sobrecarga de protocolo por requisição: uma rodada HTTP mais análise do
    protocolo de linha para InfluxDB, um ciclo completo cliente-servidor PostgreSQL mais flush
    do WAL para TimescaleDB e uma única chamada Thrift RPC para IoTDB.

  \item[\texttt{batch-large}] \texttt{BATCH\_SIZE=1000}: 1.000 pontos de dados por requisição.
    Testa a eficiência de carregamento em lote e revela o teto prático da arquitetura de lote de
    cada banco de dados. A razão entre a vazão de \texttt{batch-small} e \texttt{batch-large}
    quantifica diretamente quanto a sobrecarga por requisição domina sobre o custo de
    transferência de dados.
\end{description}

Cada teste de escrita limpa e repovooa o banco de dados antes de ser executado, garantindo
medições independentes e sem histórico.

\subsubsection{Testes de Leitura}

Cinco testes de leitura cobrem os padrões de consulta mais comuns em pipelines de análise IoT:

\begin{description}
  \item[\texttt{latest-point}] Recupera o valor mais recente para cada sensor. Esta é a consulta
    fundamental de monitoramento em tempo real: um painel ou sistema de alertas que consulta
    continuamente o estado atual do dispositivo.

  \item[\texttt{downsample}] Executa consultas \texttt{GROUP~BY} de janela de tempo calculando
    médias em janelas de tempo fixas. Modela o caso de uso padrão de agregação em painel, onde
    leituras brutas de sensores são consolidadas em resumos de nível de minuto ou hora.

  \item[\texttt{range-query}] Recupera todos os pontos de dados dentro de uma janela de tempo
    especificada, sem filtragem por valor. Modela exportação de dados históricos, pipelines de
    reprodução e ingestão para retraining de modelos.

  \item[\texttt{value-filter}] Aplica predicados de valor como \texttt{temperature~>~85.0},
    incluindo agregações sobre intervalos filtrados por valor. Modela alertas por limite:
    ``encontrar todas as leituras de sensores que excedem o nível de alarme na última hora.''
    O filtro de valor é o tipo de consulta mais discriminatório arquiteturalmente neste benchmark
    porque não pode ser respondido a partir de índices baseados em tempo: cada banco de dados
    deve manter um índice secundário ou varrer todos os valores armazenados.

  \item[\texttt{read}] Executa todos os dez tipos de consulta suportados pelo iot-benchmark
    com peso igual, fornecendo uma linha de base abrangente de carga de trabalho mista.
\end{description}

Todos os testes de leitura operam no conjunto de dados carregado pelo teste de escrita de
linha de base anterior.

\subsection{Parâmetros de Escala}
\label{sec:scale-params}

Três escalas de volume de dados são definidas para expor como as características de desempenho
evoluem com o tamanho do conjunto de dados. A Tabela~\ref{tab:scales} resume os parâmetros.

\begin{table*}[t]
  \centering
  \caption{Parâmetros de escala do benchmark. Todas as escalas usam um tamanho de lote padrão
    de 100 linhas por requisição. A escala grande contém aproximadamente cem vezes mais dados
  que a pequena.}
  \label{tab:scales}
  \begin{tabular}{lrrrrr}
    \hline
    \textbf{Escala} & \textbf{Clientes} & \textbf{Dispositivos} & \textbf{Sensores} & \textbf{Laços} & \textbf{Duração} \\
    \hline
    Pequena &  5 & 10 & 10 & 1.000 & ${\sim}$5 min \\
    Média   &  5 & 10 & 10 & 5.000 & ${\sim}$30 min \\
    Grande  & 10 & 50 & 20 & 5.000 & ${\sim}$2 h \\
    \hline
  \end{tabular}
\end{table*}

A escala grande contém aproximadamente 20$\times$ mais dados que a média e cem vezes mais
que a pequena, impulsionada tanto por uma contagem de laços maior quanto por uma topologia
expandida (10 clientes $\times$ 50 dispositivos $\times$ 20 sensores versus
5~$\times$~10~$\times$~10). A medição em múltiplas escalas é essencial porque várias
características fundamentais (saturação de compilação JIT da JVM, esgotamento de buffer,
padrões de paralisação de checkpoint e escalonamento de varredura de índice) só se
manifestam ou se tornam dominantes em volumes de dados mais elevados. Um benchmark de escala
única arrisca tirar conclusões que se aplicam apenas a um intervalo estreito de dados.

\subsection{Infraestrutura e Reprodutibilidade}
\label{sec:infrastructure}

A reprodutibilidade neste projeto repousa em três componentes.

\textbf{Docker Compose}~\cite{docker} gerencia os contêineres de banco de dados com versões
de imagem fixadas: \texttt{apache/iotdb:1.3.0-standalone}, \texttt{influxdb:2.7-alpine} e
\texttt{timescale/timescaledb:latest-pg15}. Os contêineres são iniciados e parados
\textit{sequencialmente}: apenas um banco de dados está ativo a qualquer momento. Esse design,
introduzido após as execuções iniciais revelarem artefatos de contenção de recursos entre bancos
de dados, dá a cada banco de dados acesso exclusivo à CPU, largura de banda de memória e E/S
do host.

\textbf{Nix flake}~\cite{nixos} é uma ferramenta opcional que fornece um ambiente de
desenvolvimento declarativo com arquivo de lock fixado. Executar \texttt{nix~develop} inicia
um shell com versões exatamente especificadas de Java~17, Maven, Python~3 e ferramentas Docker.
Executar \texttt{nix~run~.\#build} compila todos os três clientes de benchmark em uma única
etapa. Toda dependência transitiva de construção é registrada no arquivo de lock do flake,
eliminando as falhas de variação de ambiente que tipicamente prejudicam replicações de benchmark.
Usuários sem Nix podem replicar os experimentos instalando manualmente as dependências listadas
(Java~17, Maven, Python~3 e Docker) em qualquer host Linux compatível.

\textbf{CLI Python} (\texttt{benchmark.py}) orquestra o fluxo de trabalho completo do benchmark
programaticamente. Ele inicia o contêiner alvo, escreve os parâmetros corretos no arquivo
\texttt{conf/config.properties} do cliente de benchmark, invoca o cliente, consulta
\texttt{docker~stats} em intervalos de um segundo para coletar amostras de CPU e memória, e
acrescenta uma linha de resultado estruturada à tabela de saída após cada teste. Uma execução
completa de nove testes e três bancos de dados em qualquer escala se resume a:

\begin{verbatim}
  python3 benchmark.py --scale large
\end{verbatim}

Esse nível de automação torna os resultados rastreáveis a um ambiente de execução específico e
versionado. Qualquer pesquisador com o repositório e o Docker disponíveis pode reproduzir o
conjunto de dados completo do zero.

\subsection{Pacote de Replicação}
\label{sec:replication-pt}

A infraestrutura completa de benchmark está publicamente disponível em
\url{https://github.com/LuizFerK/iot-benchrunner}. Pesquisadores podem
reproduzir o ambiente experimental completo em qualquer host Linux
utilizando o toolchain Nix ou pacotes de sistema convencionais.

\textbf{Requisitos.} Java~17, Maven, Python~3 e Docker com Docker Compose
são necessários. Usuários Nix obtêm todas as dependências automaticamente
via o flake fornecido; usuários não-Nix devem instalar as mesmas versões
manualmente.

\textbf{Configuração.} Clone o repositório incluindo o submódulo
\texttt{iot-benchmark}, em seguida compile os clientes Java uma única vez:

\begin{verbatim}
git clone --recurse-submodules \
  https://github.com/LuizFerK/iot-benchrunner
cd iot-benchrunner

# Nix (recomendado)
nix run .#build

# Sem Nix
cd iot-benchmark
mvn -pl influxdb-2.0,timescaledb,iotdb-1.3 \
  package -DskipTests -q
cd ..
\end{verbatim}

\textbf{Execução.} Uma rodada completa com nove testes e três bancos de
dados em qualquer escala resume-se a um único comando. Usuários Nix devem
primeiro entrar no ambiente de desenvolvimento com \texttt{nix~develop},
que fornece todas as dependências de execução.

\begin{verbatim}
python3 benchmark.py --scale small   # ~5 min
python3 benchmark.py --scale medium  # ~30 min
python3 benchmark.py --scale large   # ~2 h
\end{verbatim}

\textbf{Geração de gráficos.} Após o benchmark produzir os CSVs de
resultados no diretório \texttt{results/}, o script de visualização
regenera todas as figuras utilizadas neste trabalho:

\begin{verbatim}
python3 charts.py --source all \
  --language pt-br
\end{verbatim}

O parâmetro \texttt{--source} aceita escalas individuais
(\texttt{small}, \texttt{medium}, \texttt{large}), variantes de
reexecução e \texttt{comparison} para gráficos de sobreposição entre
escalas. Os arquivos PNG de saída são gravados no diretório
\texttt{charts/}.

\subsection{Métricas}
\label{sec:metrics}

A Tabela~\ref{tab:metrics} descreve as métricas de desempenho coletadas para cada execução.

\begin{table*}[t]
  \centering
  \caption{Métricas de desempenho coletadas por execução de benchmark.}
  \label{tab:metrics}
  \begin{tabular}{ll}
    \hline
    \textbf{Métrica} & \textbf{Fonte} \\
    \hline
    Vazão (pts/s)            & Matriz de Resultado do iot-benchmark \\
    Latência média (ms)      & Matriz de Latência do iot-benchmark \\
    Latência P99 (ms)        & Matriz de Latência do iot-benchmark \\
    CPU média/pico (\%)      & \texttt{docker stats} (intervalo de 1\,s) \\
    Memória média/pico (MB)  & RSS do \texttt{docker stats} \\
    Tempo de parede (s)      & Duração da execução \\
    \hline
  \end{tabular}
\end{table*}

Para os testes de escrita, a vazão mede pontos de dados ingeridos por segundo. Para os testes
de leitura, a vazão soma as contagens de pontos retornados por todos os tipos de consulta ativos.
Como diferentes tipos de consulta retornam quantidades fundamentalmente diferentes de dados por
chamada: uma consulta de intervalo retorna centenas de pontos, uma consulta de último valor
retorna exatamente um. A \textit{latência} é portanto usada como a métrica primária de comparação de
leitura em toda a Seção~\ref{sec:experiments}.

A latência P99 captura o comportamento de cauda: 99\% das operações foram concluídas mais
rápido que esse valor. A diferença entre a média e o P99 revela a consistência do tempo de
resposta. Uma grande diferença sinaliza outliers severos ocasionais que produziriam jitter
visível em aplicações IoT de monitoramento ao vivo, mesmo quando a latência média parece
aceitável.

\subsection{Ajuste de Configuração}
\label{sec:tuning}

Após completar as execuções iniciais de benchmark, duas melhorias metodológicas foram aplicadas
antes de reexecutar todas as escalas:

\begin{enumerate}
  \item \textbf{Ciclo de vida sequencial de contêineres.} A configuração original iniciava todos
    os três contêineres de banco de dados simultaneamente. Embora cada teste fosse executado em
    um banco de dados por vez, os contêineres ociosos competiam pelo cache de páginas do SO e
    largura de banda de E/S~\citep{shapira2016pitfalls}. O TimescaleDB acumulou vários gigabytes de cache de páginas durante
    as sessões de escrita do IoTDB, produzindo medições de memória profundamente enganosas. A JVM
    do IoTDB pré-alocou seu heap máximo configurado durante o tempo ocioso pela mesma razão.
    A configuração corrigida inicia cada contêiner imediatamente antes de seus testes e o para
    imediatamente após.

  \item \textbf{Ajuste no nível do banco de dados.} O TimescaleDB recebeu uma configuração
    pgtune~\cite{pgtune} para PostgreSQL~15 direcionada ao hardware desta máquina:
    \texttt{shared\_buffers=8GB}, \texttt{work\_mem=27413kB},
    \texttt{maintenance\_work\_mem=2GB}, \texttt{effective\_io\_concurrency=1000} e
    \texttt{checkpoint\_completion\_target=0.9}. Essas configurações são a linha de base padrão
    para qualquer implantação PostgreSQL em produção nesta classe de hardware; os padrões são
    ajustados para ambientes com RAM mínima e produzem desempenho artificialmente ruim em
    máquinas capazes. O teto do cache de escrita TSM do InfluxDB foi elevado do padrão de 1\,GB
    para 8\,GB para corresponder à RAM disponível. Um limite de 28\,GB de memória foi aplicado a
    cada contêiner.
\end{enumerate}

Tanto os resultados originais quanto os ajustados são reportados e analisados na escala pequena,
onde o contraste é mais instrutivo. Os resultados de escala média e grande utilizam a metodologia
ajustada ao longo. A comparação entre as execuções original e ajustada serve a um duplo
propósito: quantifica a magnitude dos artefatos de contenção e valida quais características de
desempenho são propriedades arquiteturais genuínas de cada banco de dados.
